\documentclass[10pt]{article}
\usepackage[a4paper,margin=21mm]{geometry}
\usepackage{amsmath,amssymb,amsthm,booktabs,array}
\usepackage[T1]{fontenc}
\usepackage{lmodern,microtype}
\usepackage[colorlinks=true,linkcolor=blue,urlcolor=blue,citecolor=blue]{hyperref}
\usepackage{fancyhdr}
\pagestyle{fancy}\fancyhf{}
\fancyhead[L]{Conjecture 00000000257}\fancyhead[R]{C0ldSmi1e}
\fancyfoot[C]{\thepage}\setlength{\headheight}{14pt}
\setlength{\parindent}{0pt}\setlength{\parskip}{4pt}\setlength{\emergencystretch}{2em}
\newcommand{\lean}[1]{\texttt{\detokenize{#1}}}
\newcommand{\Q}{\mathbb{Q}}\newcommand{\N}{\mathbb{N}}
\newcommand{\R}{\mathbb{R}}\newcommand{\C}{\mathbb{C}}
\newtheorem*{theorem}{Theorem}
\begin{document}
\begin{center}
{\LARGE A fixed field along an unbounded Pell family}\\[5pt]
{\large Disproof of conjecture 00000000257}\\[5pt]
C0ldSmi1e \quad --- \quad 5 October 2026
\end{center}

The parameter $n$ can grow without bound while
$\Q(\sqrt{-(n^2+1)})$ remains the same imaginary quadratic field.
Its fixed finite class number therefore cannot satisfy any positive
uniform logarithmic lower bound. The entire argument below is formalized
in Lean using Mathlib's actual number-field class number.

\section{Exact statement and interpretation}
The English source reads:
\begin{quote}\small
Definition: $h(-d)$ is the class number of the imaginary quadratic field
$\Q(\sqrt{-d})$. Conjecture: $h(-(n^2+1))\geq c\cdot\log n$ holds for an explicit constant
$c>0$ and all $n\geq2$.
\end{quote}
The exact English and Chinese source bytes are supplied in
\lean{conjecture.md}. The two versions agree. We use natural-number $n$
(the customary integer interpretation) and the natural logarithm. We prove
that \emph{no} positive real constant works, which also rules out an explicit
one. The source imposes no squarefree condition and defines a field class
number, not a quadratic-order class number.

\section{Unbounded negative Pell solutions}
For every $k\in\N$ there exist $n,m\in\N$ such that
\[
 k<n,\qquad m>0,\qquad n^2+1=2m^2.                 \tag{1}
\]
For $k=0$, use $(n,m)=(1,1)$. Given a solution for $k$, set
\[
 n'=3n+4m,\qquad m'=2n+3m.
\]
Then $m'>0$ and $n'>k+1$, since $n>k$ and $m\geq1$. Direct expansion gives
\[
 (3n+4m)^2-2(2n+3m)^2=n^2-2m^2=-1,
\]
so $(n',m')$ is a solution for $k+1$. This proves (1) by induction, including
unboundedness; no numerical search or unproved Pell theorem is needed.

\section{The field and its class number stay fixed}
Let $K_d=\Q(i\sqrt d)\subseteq\C$ for $d\in\N$. For $m\in\N$ with $m>0$,
\[
 i\sqrt{dm^2}=m\,i\sqrt d,\qquad K_{dm^2}=K_d.    \tag{2}
\]
Both containments follow because multiplication and division by the nonzero
rational number $m$ stay inside the field. In particular, every solution in
(1) gives $K_{n^2+1}=K_2$. Write $H=h(-2)$, the finite class number of $K_2$.
Thus
\[
 h(-(n^2+1))=H                                    \tag{3}
\]
throughout this unbounded family. We do not need to calculate $H$.
The radicand is not being treated as a field discriminant: (2) is an equality
of actual embedded fields and hence of their full rings of integers.

\section{Contradiction for every positive constant}
Fix $c>0$. Choose $k\in\N$ with $k>\exp(H/c)$. Apply (1) with $k+2$ to get
$n>k+2$ and $m>0$. Then $n\geq2$, (3) holds, and
\[
 n>\exp(H/c)\quad\Longrightarrow\quad H/c<\log n
 \quad\Longrightarrow\quad h(-(n^2+1))=H<c\log n.
\]
This contradicts the proposed lower bound for that $n$.
\begin{theorem}
For every $c\in\R$ with $c>0$, there is $n\in\N$, $n\geq2$, such that
$h(-(n^2+1))<c\log n$. Consequently conjecture 00000000257 is false.
\end{theorem}

\newpage
\section{Faithful formal objects and theorem coverage}
All named declarations below belong to namespace \lean{Conjecture257}.
\begin{center}\small
\begin{tabular}{@{}>{\raggedright\arraybackslash}p{.44\textwidth}>{\raggedright\arraybackslash}p{.52\textwidth}@{}}
\toprule
Lean declaration & Mathematical content\\
\midrule
\lean{imaginaryRoot}, \lean{quadraticField}
& $i\sqrt d$ and the literal $\Q$-intermediate field of $\C$ generated by it.\\[3pt]
\lean{imaginaryRoot_sq},\newline \lean{imaginaryRoot_integral}
& The generator squares to $-d$ and is algebraic over $\Q$.\\[3pt]
\lean{quadraticFieldNumberField}
& A genuine characteristic-zero finite extension of $\Q$.\\[3pt]
\lean{imaginaryRoot_im_pos},\newline \lean{definingPolynomial_irreducible},\newline \lean{quadraticField_finrank}
& For $d>0$, the generator has positive imaginary part and $X^2+d$ is irreducible over $\Q$; the field has degree exactly two.\\[3pt]
\lean{imaginaryClassNumber}
& Mathlib's \lean{NumberField.classNumber} of this actual field.\\[3pt]
\lean{quadraticField_mul_square},\newline \lean{imaginaryClassNumber_mul_square}
& Literal equality of embedded fields, then equality of their standard class numbers, for $m\ne0$.\\[3pt]
\lean{pell_unbounded}
& The quantified unboundedness statement (1), proved by induction.\\[3pt]
\lean{counterexample_for_every_constant}
& The stronger counterexample theorem for every real $c>0$.\\[3pt]
\lean{Conjecture}, \lean{conjecture_false}
& The exact positive-constant existence claim and its closed negation.\\
\bottomrule
\end{tabular}
\end{center}

Mathlib defines the full ring of integers as the integral closure of
$\mathbb Z$ in the number field. Its class number is the cardinality of the
ideal class group, with finiteness proved in the library \cite{classnumber}.
This is not an infinite cardinal coerced to zero, an arbitrary function of
the radicand, or the class number of a varying nonmaximal order.
The intermediate-field model uses the ordinary field generated by the
specified root \cite{adjoin}. All field-identity and adequacy statements above
are proved in the submitted Lean source; none is assumed as an extra axiom.

The final proposition is exactly
\[
 \neg\exists c\in\R,\quad 0<c\ \land\
 \forall n\in\N,\quad 2\leq n\ \Longrightarrow\
 c\log n\leq h(-(n^2+1)).
\]
A larger domain containing the positive integers cannot rescue the uniform
claim. Changing to any fixed logarithm base greater than one only rescales
the positive constant, so it does not change the disproof.

\section{Verification and reproduction}
The complete project uses Lean 4.19.0 and Mathlib v4.19.0 with all nine
package revisions pinned. Separate clean project builds and strict replays
of all three source modules passed. The audit covers all 17 handwritten
declarations (4 definitions, 12 theorems, 1 instance) and all 5 generated
logical declarations. All 22 use only \lean{propext}, \lean{Classical.choice},
and \lean{Quot.sound}; their transitive logical dependencies contain no unsafe,
partial, or missing declaration. No admitted proof, \lean{native_decide}, or
custom axiom is used.

No auxiliary numerical computation supports the mathematics: the
unboundedness and field/class-number identities are formal theorems.
Verification scripts only execute or inspect the Lean project. Exact source
hashes, command logs, complete dependency and axiom inventories, and the
independent local semantic review are supplied in \lean{verification/}.
These are local validation records; maintainer acceptance remains pending.

\newpage
\section{Commands and provenance}
With the pinned Lean toolchain installed, run from the submission folder:
\begin{verbatim}
cd lean
lake exe cache get
lake build
lake env lean -DwarningAsError=true Conjecture257.lean
lake env lean -DwarningAsError=true Audit257.lean
lake env lean -DwarningAsError=true Check.lean
\end{verbatim}
Keep \lean{lake-manifest.json}; it pins Mathlib to commit
\begin{center}\small\lean{c44e0c8ee63ca166450922a373c7409c5d26b00b}.\end{center}
The initial dependency/cache retrieval requires network access. The local
verification used fresh project output directories with pinned standard
library and Mathlib dependency artifacts; it did not rebuild all of Mathlib
from source. The compiler identity and dependency revisions and tracked
sources were checked before and after execution.

\lean{Audit257.lean} prints axiom dependencies of every handwritten
declaration and instance. \lean{Check.lean} also displays all definitions and
theorem types. The separate execution inspector in \lean{verification/}
inventories all declarations originating in the submitted modules, including
generated ones, and traverses logical dependencies. Its metaprogramming
reads existing declarations and supplies no assumption or proof to the
mathematical project. \lean{README.md} gives the full inspector command with
machine-specific paths supplied explicitly.

The mathematical argument was developed independently from the supplied
bilingual statement and pinned Mathlib definitions. No previous competition
solution or problem-specific external argument was used. No external
class-number computation, approximation to logarithms, or finite numerical
prefix is substituted for the quantified theorem. Earlier failed compiler
attempts are distinguished from the successful frozen verification records.

The standalone source \lean{report.tex} compiles to the supplied
\lean{report.pdf}. The final PDF is visually checked page by page, and its
source/output hashes are recorded with the verification evidence.

\begin{thebibliography}{9}\small\raggedright
\bibitem{classnumber} Mathlib contributors,
\emph{Class numbers of number fields},
\lean{Mathlib/NumberTheory/NumberField/ClassNumber.lean},
at the pinned v4.19.0 revision.
\href{https://github.com/leanprover-community/mathlib4/blob/c44e0c8ee63ca166450922a373c7409c5d26b00b/Mathlib/NumberTheory/NumberField/ClassNumber.lean}{Source: class-number definition and proved finiteness}.
\bibitem{adjoin} Mathlib contributors,
\emph{Adjoining elements to intermediate fields},
\lean{Mathlib/FieldTheory/IntermediateField/Adjoin/Defs.lean},
at the same pinned revision.
\href{https://github.com/leanprover-community/mathlib4/blob/c44e0c8ee63ca166450922a373c7409c5d26b00b/Mathlib/FieldTheory/IntermediateField/Adjoin/Defs.lean}{Source: the generated intermediate field}.
\end{thebibliography}
\end{document}
